\chapter{Value, Scale, and Overhead}

Chapter 3 left the reading able to travel and answer for itself, with one thing it still lacked: a number. Chapter 4 is where the number arrives. The dial gets a scale --- units --- and with the scale the reading becomes a magnitude: not just that the car is going, but how much, a size that can be said aloud and a size that can be set against another. That is what this chapter adds.

A number is dangerous in a way a route is not. You can read sixty and act on it without walking back through the wheel-turns that made it, and that shortcut is how a reading forgets where it came from. So before the dial is allowed its number, the chapter makes the number answer for its origin: it must still point back to the route that produced it, to the act that carried it, and stay reusable without falling into amnesia. You may read sixty and use it; you may not read sixty and forget it was made.

Once a reading has a size it can face another --- sixty against thirty, this stretch against that. That facing is what the scale is for: a shared rule of units that lets two sizes be compared without either losing the route that made it. But the facing is not free. To hold sixty against thirty honestly, the apparatus has to keep the comparison rule --- the units, the limits --- alive and callable, and that kept obligation is a load it now carries for having made two readings comparable.

The scale has to live somewhere, and the dial's face is finite: a bounded arc, nought to some top mark, every reading standing at a place on that finite face. The scale is not spread across an endless line; it is laid on a finite frame, and that finite frame is the first place the chapter has where readings can stand and be found. Keeping it usable is its own standing cost --- the markings legible, the calibration true, the face lit enough to read --- and the chapter names that upkeep overhead: not waste, but the price of an instrument that can still answer when challenged.

And here the dial makes a claim it can almost not keep. It says: this much, sixty exactly, here, at this one place on the face, this instant, and nothing has yet come to doubt it. But you and I, watching the needle, know better --- press the scale finer to pin the size exactly and the place it sits, on a face with only so many slots, begins to swim; hold the place dead still and the last figure of the size will not stop trembling. The dial offers a sharp number at a sharp place as though it owed nothing for the pairing. It owes something. The chapter does not say what yet --- it only lets the dial make the claim, and lets us notice the claim is a little too good to be true.

\section{Origin, Enactment, and Abstraction}

Chapter 3 leaves a bounded observation able to travel: it carries a receipt, a local scene, and enough custody for another position to challenge it. That achievement does not yet say what the observation contributes --- how one carried observation should face another, how much any difference should count, or what burden the comparison places on the apparatus. Chapter 4 begins with a narrower question before those larger ones. How can a carried observation claim an origin, pass through enactment, and become available for disciplined reuse, without losing the receipt that made it answerable?

Origin does not mean a hidden birthplace. It names the point at which a procedure claims authority over what follows, and the claim must carry a receipt, because mere ancestry cannot govern a later passage. A mark may come from many conditions, but only an origin claim fixes the beginning whose custody the procedure will answer for. The apparatus therefore treats origin as a responsibility rather than an excuse: it asks what the claim permits, what it excludes, and which later acts must still answer to it. An origin that cannot return as a receipt gives the apparatus only an assertion; an origin that can return as a receipt gives the passage a place from which challenge can begin.

The origin claim also needs an inscription. The inscription need not explain the whole world; it only needs to carry the route grammar the later act will follow. The apparatus tanges that grammar into a carried instruction mark, which names which distinction begins the passage, which alternatives matter, and which transition may follow another, narrowing authority into a usable form. Without such a mark the origin claim stays too loose --- it can bless anything after the fact --- while with the mark the apparatus gains a handle that can travel through the next act and still point back to the authority it carried.

Enactment brings that authority to consequence. The apparatus does not honor an origin claim merely by preserving it; it must carry the claim through a finite route and return an outcome another position can inspect. Enactment therefore differs from announcement: announcement says a route ought to matter, while enactment lets the route do the work it claimed it could do, tanging the origin claim into a consequence. That distinction protects the receipt. If the origin claim never meets consequence, the receipt can only witness an intention; once the apparatus carries out the route, the receipt witnesses a passage --- this claim began here, this instruction traveled here, and this consequence came back under that custody.

This carrying-out also gives comparison its first grip. A consequence does not float free of the route that produced it; it returns with the shape of the passage still attached, so the apparatus can ask whether the returned outcome follows the instruction mark, whether the route left residue at a boundary, and whether another route could face the same demand. Enactment turns an origin claim into something accountable by placing the claim where it can succeed, fail, or expose a further obligation. Later chapters will ask more of this accountability; this section needs only the first discipline --- an origin earns weight by entering a carried procedure and returning with consequences.

Disciplined reuse begins after a passage works. The apparatus may want to use the successful passage again, or let another route depend on it, without dragging every detail of the earlier passage into the new scene --- and that desire creates danger, because reuse can become amnesia if it lets a solved route impersonate a primitive thing with no custody behind it. The apparatus therefore treats the later shortcut as reuse under discipline: a solved route may funge into later work only when its receipt preserves enough of the origin, inscription, and enactment for a later challenge to find its way back. Writing the answerable claim as its three parts,
\[ c = (\,o,\ e,\ a\,), \qquad a \longrightarrow (o, e), \]
with $o$ the origin, $e$ the enactment, and $a$ the abstraction that lets the claim stand for later work, the arrow is the discipline: the abstraction is reusable only while it still traces back to the origin and enactment that produced it. The contribution is finite and produced --- not a function summoned from nowhere, not a limit, not a real-valued measure.

This discipline makes reuse productive rather than magical. A reusable passage does not ask the next route to relive every prior event; it asks the next route to trust a receipt that still knows how to answer. The apparatus can then shorten future work without hiding the work that made the shortening honest: this passage has already carried its claim to consequence, so the returned shape may funge forward and be used here, provided the receipt stays near enough for a challenge to call the earlier custody into view. Reuse saves effort only because custody remains available --- when custody disappears, reuse no longer carries a solved route but an unanswerable shortcut.

Same-after-procedure gives this reuse its sharper form. Two marks need not start identical to face one another after the apparatus has carried them through a declared passage; the relevant sameness comes after the work. It means the route has explained which differences mattered, which fell away, and which receipt now lets the marks stand together. The procedure does not erase history --- it tells the observer what can count as the same once the finite hand of the apparatus has finished its declared acts.

That finite hand matters, because counting begins with discrete care before it becomes large. A carried procedure touches one distinction, then another, then another, and must know when it has made a step, when it has returned a consequence, and when it may repeat the passage without changing the custody that authorizes it. This counted hand gives Chapter 4 its first pressure: value cannot arrive as a decorative name for a successful observation, and magnitude, scale, and load cannot arrive as loose intensities. Each must answer to origin, enactment, disciplined reuse, and same-after-procedure, and the next section begins when the carried receipt asks what it contributes, how much that contribution counts, how it changes under comparison, and what burden the apparatus must continue to carry. A contribution is worth nothing the apparatus cannot still trace back to the act that produced it.

\section{Value, Magnitude, Scale, and Load}

Disciplined reuse changes the question from passage to contribution. A procedure has already claimed an origin, entered enactment, and returned with a receipt another position can challenge, so the apparatus can reuse that passage without repeating every event that made it honest. Yet reuse alone says only that the passage may travel; it does not say what the passage gives to the next act, how much of that giving has gathered, how another act may compare with it, or what burden remains when the comparison must still answer to the receipt. Value begins at that pressure. It names what a carried passage tanges into later work under custody.

Value therefore does not name a fee, an appetite, a prize, applause, or a decree; those names attach themselves to a passage from outside and flatter it after the fact, and the apparatus needs something stricter. Value names accountable contribution under a receipt: the part of the passage another act may take up because the earlier act returned with enough custody to answer for it. A receipt does not merely remember that something happened --- it licenses a later reckoning, letting the apparatus ask what this passage added to the possible work of the next, and whether that addition can face a challenge without borrowing a hidden account.

Contribution is the right word, because a passage does not carry value by wearing a label; it carries value when it changes what later work may do while remaining answerable for the change. A successful route may contribute a distinction, a confirmed custody, a narrowed alternative, a preserved residue, or a reusable answer to a local demand, and each case has a before and an after: before the route, the next act lacks some licensed way to proceed; after it, the next act may proceed through the returned shape. Value lives in that before/after difference under receipt --- the accountable difference the apparatus can cite when it lets the later act inherit the earlier work. In the notation of the previous section, value is what the contribution $c = (o, e, a)$ gives once its abstraction $a$ stands for later work: the difference $a$ licenses, kept tethered to the origin and enactment that produced it.

Value also needs answerable comparison. A lone contribution can tempt the apparatus into praise, but praise does not teach it how to face a second contribution; the receipt must support a comparison rule. One passage may contribute a sharper distinction, another a broader custody, a third a route that saves later acts from reopening the origin claim, and the apparatus cannot collapse those differences into a single mood. The comparison stays answerable only because it does not pretend that all contributions share one face: the apparatus must declare the aspect under which it compares them. Asked how much custody each passage leaves available, it owes a custody comparison; asked how much later work each can shorten, a shortening comparison; asked how much residue each preserves, a residue comparison. In every case value stays tied to an answerable reckoning rather than drifting into approval.

Magnitude enters when contribution has gathered enough extent for the apparatus to ask how much lies under custody. It does not mean raw size, which can count bulk without asking whether any part passed through the receipt. Magnitude names accumulated extent --- gathered contribution the apparatus admits because the gathering kept custody visible --- and the apparatus tanges that gathering into a finite count:
\[ \mathrm{mag}(v) = n \in \mathbb{N}, \]
where $v$ is the value and $n$ the finite extent it has gathered. The question turns from what a passage contributed to how much accountable contribution a passage, or a chain of passages, has gathered for later use, and the answer must still belong to the route grammar that made the contribution countable. The magnitude is a count under custody, not a real-valued measure of bulk.

Accumulated extent requires a disciplined gathering process. The apparatus cannot heap contributions together merely because they appear near one another; it must know which contribution may follow which, which receipt covers the gathering, and which comparison rule permits the collected result to stand as one extent. If a passage contributes a distinction and another a custody check, the apparatus must decide whether the present reckoning can gather them as one admitted accumulation. The gathering carries that decision: these contributions now stand together under this receipt, for this comparison, with these limits on later use.

Magnitude therefore gives value a public how-much without freeing it from custody. The apparatus may say that one route gathered more accountable contribution than another only after it names the kind of extent at stake --- more alternatives narrowed, more residue preserved, more later acts licensed, or more custody kept intact across a longer passage. The bare word more has no honest force until that kind is declared. The guard matters most when a passage travels into a new position: a contribution may look large in the scene that produced it and small in a scene that asks a different question, or modest near its origin and decisive once later work needs exactly the custody it preserves. Magnitude does not rescue the apparatus with a view from nowhere; it records how much lies under custody for the declared reckoning. The how-much belongs to the comparison rule, the route, and the receipt together, and if any one drops away, magnitude loses its claim.

Scale begins when the apparatus carries accumulated extent through compatible comparison. It is not a yardstick laid over every passage, and not a loose intensity that makes one contribution feel stronger than another. Scale names extent turned through comparison: the apparatus holds an accumulated extent in one direction of reckoning, then asks how that extent can funge forward to face a second direction, a later moment, or another compatible demand. Writing the turn as a comparison under a declared scale $\sigma$,
\[ v \preceq_{\sigma} v', \]
one value is counted against another only along a declared, compatible route --- a finite ordering, not a metric or a norm. The turn must preserve what the receipt can still answer for; otherwise comparison becomes conversion without custody, and the apparent scale only hides a change of question.

A declared scale has two duties. First, it names which extent travels --- accumulated custody, accumulated distinction, accumulated residue, or accumulated permission for later work. Second, it names which comparison can receive that extent without breaking it, since a receipt that supports one comparison may not support another. The apparatus must not treat every extent as directly interchangeable with every other; scale protects the difference by naming the compatible turn, letting the apparatus say that this much under this receipt may face that much under that demand while keeping the declared route of comparison in view.

This is why scale belongs between magnitude and load. Magnitude gives a how-much under custody, but one that has not yet faced a second direction; scale carries that how-much into a relation where another demand can meet it --- comparing passages in sequence, two positions inside one larger custody, or a local contribution against a later obligation. In each case the apparatus must preserve two things at once: the extent that travels and the compatibility that lets it travel; if either fails, it has not scaled the contribution but merely renamed it. Scale also stops a common confusion. The apparatus gains no power by making one grand comparison swallow every local contribution; a declared scale narrows comparison as much as it enables it, naming which aspect travels, which stays behind, and which residue the comparison must leave for a future act. The scale may enlarge a contribution's reach, but it adds discipline --- a scaled contribution has more public work to do than an unscaled one, answering both for its accumulated extent and for the compatibility of the turn that carried it elsewhere.

Load names the burden that appears when a scaled contribution must continue through later work. Once the apparatus declares a scale, the comparison has not ended the obligation: the scaled contribution now has to remain carried, constrained, and answerable as later acts depend on it. Load does not name mere heaviness but carried burden --- the obligation the apparatus accepts when it lets a scaled contribution stand as a premise for future passages. Writing the burden of a comparison,
\[ \mathrm{load}(v \preceq_{\sigma} v') = \text{the receipt and route the comparison must keep available}, \]
the more a later act rests on the scaled contribution, the more the apparatus must keep that receipt able to funge back under challenge. The load is a finite obligation, not a magnitude of effort.

This burden has several faces but one rule: the apparatus must keep the contribution from escaping the comparison that gave it standing. If a later act uses a scaled custody claim, the apparatus must remember which custody traveled and which compatibility made the travel honest; if it uses a scaled residue claim, it must preserve the route by which the residue stayed inspectable; if it uses a scaled permission, it must keep the limits of that permission near the later act. Load forces the apparatus to carry not only an answer but the discipline that made the answer usable. The burden grows when a scaled contribution becomes a junction for further work: a contribution no later act uses can rest with its receipt, but one that later acts repeatedly cite becomes a place where many routes depend at once. The apparatus must then ask whether the scaled contribution can bear those routes without losing custody, tracking which later demands belong to the declared scale and which ask for a new comparison. Load keeps that distinction from vanishing, making dependence visible instead of letting later work treat the scaled contribution as a free-standing thing.

Load completes the local chain that began with disciplined reuse. Value says what contribution the receipt can make accountable; magnitude, how much accountable contribution the apparatus has gathered under custody; scale, how that accumulated extent turns through compatible comparison; load, what burden the apparatus carries when later work keeps using the scaled contribution. Each term adds an obligation rather than a decoration. All four are needed, because a reusable passage that contributes something, gathers extent, travels through comparison, and bears later dependence has become more than a solved route. It has become part of the apparatus that future routes must answer through. None of the four is a real-valued measure: the magnitude is a count, the scale a declared ordering, the load a finite obligation. They let contribution be quantified, compared, and carried without importing a continuum the apparatus could not answer for.

The section stops here because load immediately asks a new question. If the apparatus carries an answerable burden under scale, where can that burden stand? It cannot stand everywhere merely because later work wants it; it must find a support frame, a bearing order, a place where a scaled contribution can stay constrained without losing the receipt that made it public. The next section takes up that demand. For now, value, magnitude, scale, and load give the apparatus the ladder of contribution --- what the passage gives, how much has gathered, how that extent turns, and what burden the apparatus must keep in custody. A contribution becomes part of the apparatus only by accepting, at every rung, a burden it can still be made to answer for.

\section{Finite Support and the Place Load Can Stand}

Load leaves the apparatus with a question of place. A scaled contribution has already given the later passage something to use, gathered accountable extent, and turned that extent through a compatible comparison; the burden now has to stand somewhere. It cannot float through the whole possible world while still claiming a receipt, and it cannot rest on every future act and still remain inspectable. If later work depends on the scaled contribution, the apparatus must say where the dependence bears weight, where challenge may arrive, and where the receipt can still answer. Finite support names that bounded place where load can stand under custody.

This place does not shrink the burden; it gives the burden a public stance. A load that claims every place at once gives the observer no handle for challenge --- it can always shift its claim elsewhere, appeal to a wider scene, or turn a missing receipt into an unnamed promise. The apparatus refuses that flight. It asks the burden to occupy a finite support: a bearing order with enough shape for another position to ask what stands there, what comparison holds it there, and what later use may depend on it. Finitude therefore protects ambition rather than limiting it, letting the burden speak from a place where answerability can reach it.

Finite support continues the discipline the earlier steps began. Value named the accountable contribution, magnitude gathered it into extent, scale turned the extent through comparison, and load named the obligation later work carries; finite support now asks that obligation to take a stand. The apparatus chooses a bounded bearing order because comparison needs a surface of contact --- a receipt can answer only when the challenger knows where to press. Writing the scaled load of the previous section as it bears on a frame,
\[ \mathrm{load}(v \preceq_{\sigma} v') \ \text{stands on}\ F, \qquad F = \{\,p_1, \dots, p_k\,\}, \quad k < \infty, \]
the support frame $F$ is a finite set of standing places, not an infinite-dimensional space, a basis, or the measure-theoretic support of a function. The word is borrowed but the object is not: here support is a finite, inspectable arrangement of places where a burden bears, chosen so that challenge knows exactly where to arrive. Without such a frame a scaled contribution becomes a universal excuse --- it promises to bear later work but offers no determinate place where the bearing can succeed or fail.

The word support stays literal in the metaphysical sense: it names the place where the apparatus lets a burden bear weight. The apparatus tanges the burden onto that place, which may hold a simple contribution, a coupled contribution, or a carried history of contributions. Each gives the burden a form another act can inspect --- which part of the scaled contribution stands alone, which stands only with another part, and which carries prior work as a tail later work may still challenge. From this a support grammar follows. Simple standing places let the apparatus point to one accountable claim and ask whether the receipt still holds; coupled places keep two accountable directions in relation without pretending either carries the whole burden alone; carried histories preserve a sequence of standing places when later use depends on how one support follows another. The grammar does not decorate the burden --- it fixes which questions stay legal once the burden has taken its place: what stands, what stands with it, and what travels behind it.

This grammar also keeps comparison from swelling beyond its warrant. A scale may carry extent across a second demand, but the frame says where that extent actually bears: a simple standing place forbids later work from pretending a coupled relation has arrived; a coupled place forces it to preserve the relation rather than spend one side as the whole; a carried history forces it to keep the history visible for challenge. Motion is allowed without losing place. Later work need not touch the burden where the earlier act set it down --- a frame-carrying procedure funges the burden through a sequence of finite steps, each preserving the comparison route that gave it standing and keeping the receipt near enough for inspection. The procedure has a double obligation: it must let the burden continue, and it must keep the support finite. Continuation without finitude lets the burden spread until no observer can say where it bears; finitude without continuation freezes it into a local token no later work can use. Each standing point stays a place of challenge, and the route persists only because finite points keep answering.

This is where the large finite problem appears. A small frame makes one local burden answerable, but the chapter needs more: a way for a single large burden to stand without turning into an unbounded totality. The apparatus gathers many standing demands into one large frame, which may become wide, intricate, and hard to survey while remaining one finite place of responsibility --- ambition cannot escape by claiming that every possible place belongs to the claim. A large frame earns its size by keeping its parts legible. If it gathers many standing places, it must not let them melt into a blur; each place keeps a role --- some hold simple claims, some coupled claims, some the memory of how one claim reached another --- so the apparatus can ask a precise question of the whole without losing the parts that make it answerable: which standing place bears this burden, which coupled place carries this relation, which carried history explains why the burden arrived here rather than somewhere else. A burden split into scraps loses the very relations that made it one burden, so the apparatus tanges the whole claim into one bounded place of responsibility --- wide enough for a large ambition, yet bounded enough that challenge can still enter it at any standing place.

This gives the frame a public grammar. A finite standing pattern is not a diagram in the air but an accountable arrangement of places where receipts can still speak --- a beginning place, a place that turns through comparison, a place that carries the result forward. The apparatus may join these places, but it must keep the joins visible, since a hidden join would let later work borrow connection without answerability while a visible join lets a challenger follow the passage from one place to the next. A joined support passage is how continuity enters without claiming unlimited reach: one standing point calls the next, the next another, each call finite and inspectable. Continuity here means disciplined carriage, not endless spread --- the burden funges through standing points without losing the receipt that made the travel public. Skip the standing points and the apparatus has only assertion; keep them visible and it has a passage later work can challenge step by step.

The single large finite burden carries an internal ethic. It refuses to hide difficulty by splitting into unrelated scraps, and refuses to hide difficulty by declaring itself everywhere at once; it asks the apparatus to gather what belongs together and keep the gathering accountable. That ethic matters because later work will want to rest on the whole frame as one reliable support, and the frame can serve that role only while it keeps enough internal order for the observer to reopen the receipt, recover the comparison route, and test where the burden bears. The frame can grow, but growth stays a finite act: a stronger burden may need more standing points, a longer passage more joins, but it still needs standing points and joins a challenge can follow. When the burden grows, the frame grows under discipline --- more standing places, more coupled places, more carried histories, but still one finite bearing order --- and largeness never becomes a substitute for answerability. The demand for one answerable large frame is what keeps greatness from becoming evasion.

The large frame also changes the role of residue. Earlier sections let residue mark what remained after a finite decision or comparison; here residue helps test the standing place. When the burden rests in a frame, the apparatus can ask what the frame still leaves for later inspection. A burden that leaves no route back to residue has become too polished; a burden that leaves residue everywhere has refused place. Finite support gives residue a bounded place to return, letting later work ask whether the standing pattern held, whether the coupled place kept its relation, and whether the carried history still knows its route.

The observer gains a stronger role as well. Before finite support, the observer could challenge a receipt, compare contributions, or ask why a scaled burden keeps its authority; with finite support, the observer can point. The observer can say where the burden claims standing, the route by which it arrived, and the finite frame that must answer. The pointing does not make the observer sovereign --- it gives challenge a place. A metaphysics of measurement needs that place, because measurement cannot survive as a pure claim of reach; it needs a location in the grammar where receipt and burden meet.

Finite support therefore refuses two evasions at once: the empty local token that cannot carry later work, and the unbounded claim that cannot face local challenge. It asks the apparatus to carry burden through finite standing patterns, to join those patterns when continuity matters, and to keep one large finite frame when ambition grows. The result gives scope its honest form --- scope becomes accountable only when it can stand somewhere. A reach that names no place to bear it is not large but weightless; the finite frame is what gives a claim the weight of having to stand, and weight of that kind is the only ambition the apparatus can be made to answer for.

The section stops before the next burden arrives. Once load has a finite place to stand, the account is not finished, because the support frame itself now carries a pressure: it takes work to hold standing places, keep relations joined, preserve histories, and keep the large finite frame available for later challenge. Finite support has named the place where load can stand; the next section asks what pressure the act of support itself carries. A claim that can stand somewhere has earned the right to be challenged there; one that can stand nowhere has not earned the right to be believed anywhere.

\section{Overhead as Conserved Pressure}

Finite support answers where load can stand; it does not make standing free. Once the apparatus gives a burden a support frame, it must keep that frame available for challenge --- preserving the standing places, maintaining the joins, carrying the histories, and keeping the route back to receipt open. The apparatus tanges each of those duties into the account it owes, and the continuing pressure of paying them is overhead. Overhead is not a secondary annoyance after the real work has ended; it belongs to the work itself --- the pressure the apparatus carries when it lets a scaled contribution stand where later acts can depend on it.

That pressure has to remain present, because support does more than store a result --- it keeps a place answerable, saying where the burden bears, which comparisons still govern it, and how later challenge can reach the receipt that gave it standing. If the apparatus drops the pressure, the frame may still look like a frame, but it no longer carries the burden in an answerable way: the place stays named while the duty disappears. Overhead is what keeps the duty from disappearing --- the cost of making support available rather than merely declared.

For that reason overhead is not waste. Waste is expenditure the apparatus could remove without changing the accountable passage; overhead is pressure that support itself requires. A standing place must stay open to challenge, a joined place must preserve its relation, a carried history must keep enough route for later work to follow --- the apparatus may reduce confusion around those tasks, but it cannot abolish the tasks while still claiming support. Writing the overhead of a frame as the finite cost of keeping its $k$ standing places available,
\[ \mathrm{ovh}(F) = \sum_{i=1}^{k} c(p_i) \;\geq\; 0, \]
each $c(p_i)$ is the finite, non-negative cost of one place with its joins and history; the sum is finite because $k < \infty$, and it is never negative, since support cannot turn a profit --- not a continuous energy, an integral, or a thermodynamic law, but a bookkeeping balance. Nor is overhead an embarrassment or an apology: carrying it is not a confession of failure but the form answerability takes once support becomes durable.

This makes conserved support pressure a positive claim rather than a complaint. The pressure cannot simply vanish once the apparatus sets the burden down, because a later act needs more than a memory of placement --- it needs a living route of answer. So the apparatus funges the pressure forward with the burden: the pressure of place, so a challenger can find where the burden bears; the pressure of relation, so the joined parts do not drift apart; the pressure of history, so a later claim can ask how the burden arrived. These pressures belong to one act of support, and none of them is paid once and then forgotten.

The pressure also grows as the support frame grows. A single standing place asks the apparatus to keep one claim answerable; a large finite frame asks it to keep many standing places in one order, where each added place brings a duty of legibility, each join a duty of relation, and each carried history a duty of route. The apparatus may gather those duties into one frame, but gathering does not remove them --- the larger frame concentrates pressure. It lets ambition stand, and in the same act makes the cost of standing explicit.

Overhead therefore follows load without replacing it. Load names the burden later work accepts when it depends on a scaled contribution; overhead names the pressure the apparatus accepts when it keeps the place of that dependence available. Load asks what the apparatus must carry; overhead asks what the carrying requires from the support itself. A later act may treat a supported burden as usable, but the apparatus still has to preserve the grammar that made the use honest, and overhead is what guards that preservation.

This difference matters for what the next chapter will call strength. Strength does not mean the pressure has vanished; it means the apparatus can bear pressure openly without letting the burden deform its route. A weak support hides the cost of support, so later users inherit a claim without knowing what custody it demands; a stronger support names the cost and distributes the work --- keep this place legible, keep this join visible, keep this history open for return, keep this large frame open to challenge. Overhead is what gives strength its ordinary labor.

The pressure stays finite even when support lasts. It does not float above every possible future use; it follows the supported burden only along the routes the apparatus has actually made available. When later work asks for more reach, the apparatus must either extend support with new standing places or refuse the extra use --- it cannot spend yesterday's support as permission for every tomorrow. Overhead keeps that refusal honest by reminding later users that support always travels with custody. Reach that outruns its custody is not strength but a debt the apparatus has not yet been asked to pay.

The guard becomes most visible when a supported claim begins to attract an account around itself. A useful burden rarely travels alone: new acts cite it, describe it, compress it, defend it, and recruit it for further work, so a recruited account gathers --- a circulating record of why the supported burden matters, where it can travel, and who may rely on it. Such circulation can help the apparatus by keeping custody intelligible, and it can also add pressure the original frame did not yet carry. When the account funges the burden into more hands while keeping the receipt near, it carries overhead without hiding it: a good account does not replace the support frame but points back to it --- here is the place where the burden stands, here is the relation it must keep, here is the history that lets challenge return.

The same circulation can also thicken the burden. A circulating account may attach more expectation to the supported claim than the receipt can answer for, make the burden sound broader than its frame permits, or pull attention toward its own force and away from the standing place. The apparatus must account for that pressure too. It cannot solve the problem by forbidding accounts, since later work often needs a public route back to the burden; it must instead ask which pressure the account recruits and whether that pressure still answers to receipt. A recruited account has a strict boundary: it may circulate the burden but not manufacture standing, help later acts remember a route but not invent one, gather attention around a claim but not turn attention into witness. When such an account grows louder, the apparatus asks a quieter question --- does it return to the place where the burden actually stands?

This lets the chapter treat social uptake without giving it the power to decide truth. A supported burden can enter a larger order: other bearers adopt it, pass it onward, and organize further work around it. The adoption may carry real pressure --- preserving a useful route, coordinating later challenge, and keeping a large finite frame from collapsing into private custody --- but it does not create the truth of the claim. It creates another responsibility: to carry the burden without cutting it loose from receipt. An adoption order gives overhead a wider place to stand, assigning bearers to preserve the burden, repeat the route, and keep the account pointing back to the frame. The overhead is not abolished by being shared; it is conserved across the bearers that carry it,
\[ \mathrm{ovh}(F) = \sum_{j} \mathrm{ovh}(b_j), \]
so adoption redistributes the cost among named bearers $b_j$ rather than removing it. The pressure funges from one bearer to the next as the burden circulates, but the total never drops below what the frame genuinely costs to keep answerable.

The danger lies in mistaking the extension for closure. A larger order may carry a burden farther, but distance does not turn into witness, more bearers do not strengthen the receipt by their number alone, and repetition does not settle what the support frame cannot answer. The apparatus must keep the adopted burden under the discipline that governed the original standing place: where does it bear, which relation does it preserve, how can later challenge return to the receipt? If those questions fail, uptake has widened pressure while weakening answerability. An institutional bearer therefore receives a burden, not a crown --- it may keep the account alive, carry it into new scenes, and help later work find the claim, but it may not replace witness, turn adoption into evidence, or hide the residue the frame leaves. At best it helps overhead stay honest by making the continuing pressure visible: this burden still requires custody, this account still owes its route, and this adoption still waits on later challenge.

Overhead also protects the apparatus from a false purity --- the wish for a claim that stands without pressure, circulates without an account, and enters a larger order without adoption. Such a wish would erase the actual labor of answerability. The apparatus works in finite frames; it needs places, joins, histories, and bearers, and each need creates pressure. The honest question is not how to make pressure disappear but how to carry it without letting it impersonate truth, evidence, or completion. This keeps the social layer subordinate to measurement: circulation and adoption change how a burden travels after support, but they do not replace the chain of origin, enactment, contribution, magnitude, scale, load, and finite support. They add a layer of custody around that chain, and the apparatus reads them as overhead --- real pressure, real cost, real public consequence, but not the final judge of the claim. Witness remains ahead, and the supported burden still owes the route by which a later act can test it.

The section ends at that boundary. Chapter 4 has now taken a passage from origin to contribution, from contribution to load, from load to finite support, and from support to the conserved pressure that support must carry; it has also shown how circulation and adoption can recruit that pressure without closing measurement. The next question asks how one apparatus can carry these obligations together --- support pressure, circulation pressure, and adopted burden, all still answerable to receipt while truth waits beyond the present ladder. Overhead is the price of keeping a claim alive enough to be doubted; a claim that costs nothing to maintain has bought its quiet by ceasing to answer.

\section{The Ladder in One Apparatus}

The preceding sections have not assembled a list of separate topics; they have drawn one ordered apparatus under increasing obligation. Origin gives the passage a place to begin, enactment gives that beginning a carried act, and abstraction lets the carried act answer beyond its first occasion. Value names the contribution the apparatus can hold to account, magnitude gathers that contribution into extent, scale turns extent through compatible comparison, load names the burden later work accepts when it depends on that comparison, finite support gives the burden a place where it can stand, and overhead names the pressure that keeps the standing place answerable. In the chapter's own symbols the rungs run in order,
\[ c = (o,e,a) \ \longrightarrow\ \mathrm{mag}(v) \ \longrightarrow\ (v \preceq_{\sigma} v') \ \longrightarrow\ \mathrm{load} \ \longrightarrow\ F \ \longrightarrow\ \mathrm{ovh}(F), \]
each rung feeding the next inside one finite apparatus --- not a category, a functor, or an infinite tower, but a single ordered run that still answers to one receipt.

This order matters because each later term receives an earlier duty rather than replacing it. The earned ladder does not throw away origin when it reaches enactment, or discard contribution when it reaches load; it carries the earlier receipt forward and adds a new burden of custody. A contribution asks the apparatus to remember what an act gave, an extent asks it to keep that giving countable, a comparison asks it to preserve the rule by which one extent answers beside another, and a burden asks it to carry the result into later work without pretending that later use erases the path by which the result arrived.

The ladder therefore has the form of successive obligation. The apparatus tanges a contribution at the first rung and funges it upward, and a repeatable mark can enter present measurement only when origin, enactment, contribution, extent, comparison, and burden stay in one answerable order. Nothing in that chain floats free. Each stage takes custody of what came before it and sharpens the next question: what can this receipt now bear, and what new responsibility follows from that bearing? The order stays earned because the apparatus does not borrow authority from a later name; it earns each name by preserving the duty that made the name available. A later stage may speak more generally, travel farther, or serve more work, but it still owes its right to the custody it accepts.

Finite support gives this carried order its standing place. Without support, load would remain a promise to carry rather than a place where carrying can face challenge. Support does not glorify smallness or retreat from reach; it says that a burden can stand only where the apparatus can keep its relations available, showing where the burden bears, which joins hold it together, and how later work can return to the receipt. Support turns the ladder from an order of names into an order of places.

Overhead keeps those places from becoming empty signs. Once the apparatus lets a burden stand, it must preserve the pressure of that standing --- keeping the place legible, the joins visible, the history open, and the route of answer available for later challenge. Overhead does not sit outside the ladder as an extra cost added after the work is done; it belongs inside the apparatus as continuing custody. The apparatus carries not only the burden but the pressure that lets the burden remain answerable.

The carried order then enters circulation without surrendering to it. Later acts need ways to remember where a burden stands and why it matters, and a recruited account can help by pointing back to the standing place, keeping the route of answer near, and coordinating later use. But that carriage remains under receipt: it may thicken attention, widen circulation, or make a route easier to repeat, yet it may not turn attention into witness. The account serves the ladder only when it carries the burden back to the place where the apparatus can still answer.

Adopted custody adds the last of the chapter's pressures. A larger order can take up a supported burden and assign bearers to keep it available, which helps when one local act no longer suffices to preserve a route --- distributing custody, stabilizing memory, keeping later challenge from losing the path home. Yet adoption receives a burden; it does not crown it. The institutional bearer must keep the route open, not turn shared carriage into truth: adoption carries farther only by accepting the discipline of receipt.

These stages form one apparatus of carried obligations. Origin, enactment, abstraction, contribution, extent, comparison, burden, support, overhead, circulation, and adopted custody all answer to the same demand: a later act must be able to ask what the apparatus has carried and why it may carry it now. The apparatus becomes ordered when it can answer that question without collapsing one duty into another --- it must not call contribution extent, extent scale, scale burden, burden support, support witness, or adoption truth. Each duty keeps its place, and the place is what lets the next-bearing order hold.

Chapter 4 therefore closes with an answerable apparatus, not with completed truth. The ladder has earned enough order for a later challenge to address it: a beginning, a carried act, an accountable contribution, a measurable extent, a comparison, a burden, a standing place, a pressure of custody, an account under receipt, and bearers who may carry the burden farther. That is a large achievement, but not the final judge --- the apparatus can now present itself for witness, and the next chapter asks what kind of witness can test this carried order without confusing circulation, support, or adoption with truth. An apparatus that can be questioned at every rung has earned the right to be tested; it has not yet earned the right to be believed.

\section*{Bridge: The Apparatus Still Needs Witness}

Chapter 4 has earned an apparatus that can answer for what it carries; it has not earned witness, and that distinction governs everything ahead. Origin gives the passage a responsible beginning, enactment carries that beginning into consequence, and abstraction lets the consequence travel without losing its receipt. Value names accountable contribution, magnitude gathers that contribution into extent, scale turns extent through compatible comparison, load names the burden later work accepts, finite support gives that burden a place to stand, and overhead names the pressure that keeps the place available. Circulation and adopted custody carry the burden farther, but they do not change what the burden still owes.

The chapter therefore closes one question by making another unavoidable. It shows how a carried act becomes an ordered apparatus of obligations in which each later duty receives an earlier one and adds a stricter demand: a contribution must remember the act that gave it standing, an extent must keep the contribution countable, a comparison must preserve the compatibility that lets one extent face another, a burden must carry the result into later work without erasing the route by which it arrived, a support frame must give the burden a place where challenge can press, and overhead must keep that place legible. Circulation and adoption must point back to the receipt instead of pretending to replace it. The apparatus can now say what it carries, where it carries it, and which obligations keep the carrying honest.

That answer remains incomplete, because answerability does not yet equal witness. An apparatus can tange a receipt, preserve a burden, and maintain a route of challenge while still lacking the public act that tests truth. The receipt records that a passage happened under custody, the support frame fixes where the burden can stand, and the overhead account states what pressure keeps that standing available --- yet none of those duties, by itself, tells the apparatus what kind of encounter can count as witness. The apparatus has prepared the route by which a witness might challenge it, but preparation and challenge differ: a stage may carry the conditions for a test without yet performing the test. The gap is exact,
\[ \text{answerable} \;\not\Rightarrow\; \text{witness}. \]

This gap protects truth from several false closures. A burden does not become true because the apparatus can carry it, a supported claim does not become true because an account makes it easy to remember, and an adopted claim does not become true because many bearers repeat it. Circulation funges a burden into more scenes, support funges it into a standing place, and adoption funges it into a wider order, yet none of those motions funges answerability into witness. Each achievement matters, and each remains subordinate. The apparatus must therefore ask a stricter question: what kind of witness can face the carried order and make truth an obligation inside the apparatus rather than a possession outside it?

The next chapter begins from that pressure. Witness must not float above the apparatus as a private certainty, and it must not collapse into public uptake; it has to enter the same discipline Chapter 4 has earned --- receipt, place, burden, pressure, route, and challenge. Truth must become something the apparatus can answer through witness, not something a bearer owns by pronouncement. Reality must meet the apparatus without letting reach outrun receipt. Decision must emerge from a finite order that can still show why this claim, this route, and this challenge belong together. Chapter 4 cannot answer those demands without beginning a new argument, so it stops at the threshold.

Chapter 4 cannot settle reality merely by naming a place where a burden stands. Finite support gives a claim a place for challenge, but reality asks whether the apparatus can meet more than its own order. A receipt can record how a passage came to stand; it cannot, by itself, say how the standing passage meets a world that can resist it. Overhead can keep the route open, but it cannot decide what a witness will find when the route opens. Adoption can carry the burden into a wider order, but it cannot make that wider order the measure of what the burden means. Reality must therefore enter as a further discipline, not as a reward for internal coherence.

The same guard applies to decision. Chapter 4 has prepared a finite order that can ask for decision without pretending decision already occurred; it has given the next chapter a burden with a place, a route, and a pressure of custody, and it has left a demand. The apparatus must learn how a witness chooses, refuses, confirms, or redirects the carried order without letting choice become possession. Decision will have to respect the receipt and the support frame while facing a challenge the apparatus cannot answer merely by repeating its own ladder.

The bridge therefore points forward without stealing the next chapter's work. Chapter 4 has not shown what scientific truth requires, how reality enters the apparatus, how locality constrains variation, how a path carries a rule, how logic ends a route, how measured output becomes public, or how inference closes on the apparatus; it has only earned the apparatus that can receive those questions without dissolving into metaphor or social pressure. Origin, enactment, abstraction, contribution, extent, comparison, load, support, overhead, circulation, and adopted custody now stand in one answerable order, and the next chapter asks what witness can do to that order: how it tests the claim, how it protects truth from possession, and how it turns an answerable apparatus toward reality and decision. An apparatus that can answer for everything it carries has done all that custody can do; it has not yet been made to face anything it cannot.

\section*{Coda: The Ladder in Review}

Chapter 4 took the traveling reading and gave it a size --- a magnitude, on a scale, comparable to another, standing on a finite frame and kept legible at a cost. That is the chapter in review. But set down inside that same accounting, unnamed, is a fourth part of the thing the chapters have been quietly assembling.

The three parts already on the bench were all sizeless: a value that can be one thing in one place and another elsewhere, a value that can be one thing before and another after, a value read through a floor of disturbance. Now the value gets a size --- a magnitude that can be said, on a scale that lets one value face another --- and the sizes need somewhere to stand. Not an endless line, but a finite frame: the places, which an earlier chapter could not yet say formed any larger thing, now form one --- a bounded arrangement where each value sits, allowed only because the finite support this chapter built gives them somewhere to stand. That is the fourth part, and it has two faces at once: a value that now has a size, and a finite frame of places the sizes are spread across.

What the frame is for does not matter right now. Like the three before it, the fourth is set down and left unexplained; what it is for waits on the parts not yet handed over. The book keeps its one habit --- one piece per chapter, in its turn, none before the construction reaches it. Four pieces now lie in order: a value spread across places, the same value moving, the value sorted from disturbance, and now the value given a size on a bounded frame of places. A reader who has watched them go down in order may feel the shape pulling tighter, though the book still will not say its name.

For now there are four parts, and the fourth is as bare as the rest. It names no size in particular, fixes no frame in particular, and leans on nothing still to come: a value that can differ from place to place, a value that can differ from before to after, a value read through a floor of disturbance, and now a value with a size, spread over a finite frame of places. It is a fourth mark of a second kind, set down beside the others and waiting, as they wait, for what has not been handed over to give it its place.
